\documentclass[12pt,a4paper]{article}
\usepackage{ugmath}
\usepackage{float}
\usepackage{placeins}
\usepackage{array}
\usepackage[labelfont=bf,labelsep=space]{caption}
\newcommand{\paren}[1]{\left( #1 \right)}
\newcommand{\alumno}{Ricardo León Martínez}
\newcommand{\materia}{Algebra Lineal I}
\newcommand{\profesor}{Claudia Reynoso Alcántara}
\newcommand{\tarea}{Examen 2}
\newcommand{\fecha}{15/5/2026}

\begin{document}

    \begin{center}
        {\large\textbf{UNIVERSIDAD DE GUANAJUATO}}\\[0.3cm]
        {\normalsize\textbf{DIVISIÓN DE CIENCIAS NATURALES Y EXACTAS}}\\
        {\normalsize\textbf{CAMPUS GUANAJUATO}}\\[1cm]

        {\Large\textbf{\tarea\ (\materia)}}\\[1cm]
    \end{center}

    
    \begin{center}
        \textbf{Fecha:} \fecha 
    \end{center}
    

    \vspace{0.3cm}

    1. (2 puntos) Encuentra una base del siguiente subespacio vectorial de $\R^{4}$
    (justifica tu respuesta).
    \begin{equation*}
        W=\{(x_{1},x_{2},x_{3},x_{4})\in\R^{4}:x_{1}+x_{2}+x_{3}+x_{4}=0\}.
    \end{equation*}

    2. (1 punto) Si existe, describe explicitamente una transformación lineal de $\R^{3}$
    en $\R^{3}$ que tenga como imagen el subespacio vectorial generado por (1,0,1) y
    (1,1,0). Si no existe, demuestralo.

    \vspace{0.5cm}

    3. (2 puntos) Sea $n\in\N$ y $\PP_{n}$ el espacio vectorial de polinomios
    de grado menor o igual a $n$ con coeficiente en $\R$. Considera la transformación
    lineal
    \begin{align*}
        T:\PP_{3}&\to\PP\\
        f(x)&\mapsto xf^{\prime}(x)
    \end{align*}
    Considera la base ordenada $\BB=\{1,x+1,x^{2},x^{3}\}$ de $\PP_{3}$. Encuentra la matriz
    que representa a $T$ en esta base ordenada.

    \vspace{0.5cm}

    4. Sea $T:\R^{4}\to\R^{3}$ la transformación lineal dada por el producto de la matriz
    \begin{equation*}
        \begin{pmatrix}
            2 & 1 & 0 & 1\\
            0 & 1 & 0 & 1\\
            2 & 2 & 0 & 2
        \end{pmatrix}
    \end{equation*}
    encuentra (y justifica)\\
    \textit{a)} (1 punto) la nulidad de $T$ y\\
    \textit{b)} (1 punto) el rango de $T$.

    \vspace{0.5cm}

    5. Sea $V$ un espacio vectorial de dimensión finita $n$, sobre $F$. Sea $T:V\to V$
    una transformación lineal tal que $T\neq0$ y $T^{2}=0$. Demuestre que:\\
    \textit{a)} $\Im(T)\subset\Nul(T)$\\
    \textit{b)} $\text{rango}(T)\leq\frac{n}{2}$.

    \vspace{2.3cm}

    6. Sean $V,W$ y $U$ espacios vectoriales de dimensión finita, sobre un campo $F$. Sean $T:V\to W$ y $S:W\to U$ transformaciones lineales.
    Demuestre que
    \begin{equation*}
        \nul(S\circ T)=\nul(T)+\dim(\Im(T)\cap\ker(T)).
    \end{equation*}

\end{document}